\documentclass[11pt,a4paper]{article}
\usepackage[margin=1in]{geometry}
\usepackage{times}
\usepackage{booktabs}
\usepackage{amsmath}
\usepackage{graphicx}
\usepackage{multirow}
\usepackage{xcolor}
\usepackage{hyperref}
\usepackage{natbib}

\title{Exception Events Are All You Need for Verdict Prediction}

\author{Anonymous}

\date{}

\begin{document}
\maketitle

\begin{abstract}
Execution traces have shown promise for improving code understanding, yet no prior work has systematically evaluated which trace components drive verdict prediction (binary pass/fail classification). We reframe verdict prediction as fault detection and conduct a component-level ablation across four models spanning three architecture families (Qwen, DeepSeek, LLaMA) at two scales. Our key finding: exception events---compact structured error summaries---consistently achieve the highest accuracy across all models tested (99.83\% on Qwen-2.5-Coder-7B with zero variance across 3 seeds; 99.83\% on CodeLlama-13B; 99.33\% on DeepSeek-Coder-6.7B), outperforming full VOTA traces by 0.42--14.88 percentage points. We further demonstrate that this performance stems from genuine fault semantics learning rather than format-level shortcuts: removing format markers produces zero accuracy degradation (McNemar $p$=1.0, 3 seeds). A two-tier diagnostic hierarchy emerges where Tier~1 signals (exceptions, assertion-points) are individually near-perfect and sufficient, while Tier~2 signals (loop compression) are model-dependent and subsumed. These results establish that verdict-optimized trace design should extract highest-specificity diagnostic signals and discard everything else.
\end{abstract}

\section{Introduction}

Execution traces capture the dynamic behavior of programs and have been explored as training signals for code language models \cite{ni2024next,jain2024tracepile}. Prior work on trace-augmented supervised fine-tuning (SFT) has focused primarily on generative tasks---code completion, output prediction---and reported mixed results. DCSH \cite{ni2024next} tested five trace formats on generation tasks and found ``no single trace representation consistently outperforms others.'' TracePile \cite{jain2024tracepile} scaled to 2.6M samples but evaluated only generation. TRACED \cite{dinella2022traced} demonstrated traces help encoder-based classification but did not isolate which trace components matter.

A critical gap remains: \emph{what trace components drive verdict prediction} (binary pass/fail classification), and \emph{why}? We argue that verdict prediction is fundamentally a \textbf{fault detection} task, and that the most effective signal is not a rich execution narrative but rather the \textbf{exception event}---a compact, structured diagnostic moment recording error type and message.

Through systematic component ablation across four models from three architecture families (Qwen-2.5-Coder-7B/1.5B, DeepSeek-Coder-6.7B, CodeLlama-13B-Instruct) we establish three claims:

\begin{enumerate}
    \item \textbf{Exception events are the most effective trace component} for verdict prediction, consistently outperforming all alternatives across all models tested.
    \item \textbf{Information density, not gradient dilution}, is the primary bottleneck explaining why trace compression helps---even oracle loss masking cannot close the 13pp gap between full traces and targeted extraction.
    \item \textbf{A two-tier diagnostic hierarchy} exists: Tier~1 signals (exceptions, assertion-points) are individually sufficient; Tier~2 signals (loop patterns) are model-dependent and subsumed.
\end{enumerate}

We further validate that exception signal effectiveness derives from fault semantics rather than format shortcuts, through a controlled format-marker ablation yielding zero performance degradation across 3 seeds.

\section{Method}

\subsection{Task Formulation}

Given a program $P$, test input $x$, and execution trace $T(P, x)$, verdict prediction is the binary classification task: $f(T(P, x)) \rightarrow \{\texttt{pass}, \texttt{fail}\}$. We study how different representations of $T$ affect classification accuracy under LoRA SFT.

\subsection{Dataset}

We construct our dataset from HumanEval (164 problems) and MBPP-sanitized (427 problems), totaling 591 problems. Each problem is solved by 10 LLM-generated solutions, yielding 5,910 samples with execution traces collected via a Python tracer. The train/test split is 4,720/1,190 samples (problem-level split, zero train-test input overlap).

\subsection{Trace Representations}

\paragraph{Full Trace.} Complete line-by-line execution trace including variable states, control flow, and all events.

\paragraph{VOTA (Verdict-Optimized Trace Abstraction).} A compressed representation that isolates three components:
\begin{itemize}
    \item \textbf{Exception events}: Structured error summaries in the format \texttt{[ERROR] type: XxxError, message: ...}
    \item \textbf{Loop compression}: Iteration patterns condensed into behavioral summaries
    \item \textbf{Actual/Expected}: Output comparison at assertion points
\end{itemize}

\paragraph{Component Ablations.} We isolate individual VOTA components:
\begin{itemize}
    \item \texttt{exception\_only}: Only exception event data
    \item \texttt{ae\_only}: Only assertion-point (actual/expected) data
    \item \texttt{loop\_only}: Only loop compression data
    \item \texttt{no\_trace}: No trace information (instruction + code only)
\end{itemize}

\subsection{Training Setup}

All models are fine-tuned with LoRA (rank 16, $\alpha$=32) for 3 epochs using SFTTrainer. Max sequence length is 2,048 tokens. We evaluate with 3 random seeds per condition on primary models.

\section{Experiments}

\subsection{Claim 1: Exception Dominance Across Models}

\paragraph{Qwen-2.5-Coder-7B (Primary).} Table~\ref{tab:qwen_main} shows results across all trace conditions. Exception-only achieves 99.83\% with perfect seed robustness (std=0.00 across 3 seeds), outperforming full VOTA (99.41\%) by 0.42pp ($p$=0.01).

\begin{table}[t]
\centering
\small
\begin{tabular}{lcccc}
\toprule
\textbf{Condition} & \textbf{Acc.} & \textbf{ECE} & \textbf{Brier} \\
\midrule
exception\_only & 99.83\% & 0.0027 & 0.0013 \\
ae\_only & 99.92\% & 0.0004 & 0.0008 \\
VOTA (full) & 99.41\% & 0.0054 & 0.0050 \\
loop\_only & 85.63\% & 0.0831 & --- \\
full\_trace & 81.76\% & 0.0958 & --- \\
full\_trace (label-only) & 86.22\% & 0.1107 & --- \\
no\_trace & 75.13\% & 0.1715 & 0.1972 \\
\bottomrule
\end{tabular}
\caption{Verdict prediction accuracy, ECE, and Brier score on Qwen-2.5-Coder-7B (seed 42). Tier~1 signals (exception, ae) achieve near-perfect accuracy \emph{and} calibration.}
\label{tab:qwen_main}
\end{table}

\paragraph{Seed Robustness (Qwen).} Table~\ref{tab:qwen_seeds} confirms zero variance across 3 seeds, with identical error samples (id=901, 905) in all runs---establishing 99.83\% as the representation ceiling.

\begin{table}[t]
\centering
\small
\begin{tabular}{lccc}
\toprule
\textbf{Seed} & \textbf{Accuracy} & \textbf{95\% CI} & \textbf{Errors} \\
\midrule
42 & 99.83\% & [99.58, 100] & 2/1190 \\
123 & 99.83\% & [99.58, 100] & 2/1190 \\
456 & 99.83\% & [99.58, 100] & 2/1190 \\
\midrule
\textbf{Mean $\pm$ Std} & \multicolumn{3}{c}{99.83\% $\pm$ 0.00} \\
\bottomrule
\end{tabular}
\caption{Qwen exception\_only across 3 seeds. All seeds produce identical errors on the same 2 samples (GeneratorExit edge cases).}
\label{tab:qwen_seeds}
\end{table}

\paragraph{Cross-Family Validation (4 Models).} Table~\ref{tab:cross_family} demonstrates exception dominance across three architecture families. CodeLlama-13B-Instruct achieves 99.83\% on exception\_only versus 92.94\% on VOTA, confirming the pattern extends to the LLaMA family.

\begin{table}[t]
\centering
\small
\begin{tabular}{llccc}
\toprule
\textbf{Model} & \textbf{Family} & \textbf{VOTA} & \textbf{exc\_only} & \textbf{$\Delta$} \\
\midrule
Qwen-2.5-7B & Qwen & 99.41\% & 99.83\% & +0.42 \\
Qwen-2.5-1.5B & Qwen & 99.58\% & 99.83\% & +0.25 \\
DeepSeek-6.7B & DeepSeek & 84.45\% & 99.33\% & +14.88 \\
CodeLlama-13B & LLaMA & 92.94\% & 99.83\% & +6.89 \\
\bottomrule
\end{tabular}
\caption{Exception dominance across 4 models $\times$ 3 families $\times$ 2 scales. Exception\_only outperforms VOTA on every model (mean $\Delta$=+5.61pp).}
\label{tab:cross_family}
\end{table}

\paragraph{CodeLlama-13B Detail.} CodeLlama-13B-Instruct on VOTA achieves 92.94\% (CI [91.4, 94.4]), with pass\_precision=0.93, pass\_recall=0.98, fail\_precision=0.93, fail\_recall=0.80. The VOTA gap pattern is consistent: models with weaker execution reasoning show larger degradation on VOTA's Tier~2 loop component, while exception signal remains uniformly high.

\subsection{Format Independence: Ruling Out Shortcuts}

A potential concern is that the model exploits a format-level shortcut---specifically, pattern-matching the marker string \texttt{[No exception data in trace]} to predict pass. We test this by removing all format markers and retraining.

\paragraph{Setup.} We remove the \texttt{[No exception data in trace]} marker from all pass samples in the exception\_only condition, leaving them as empty or minimal traces. Models are retrained from scratch across 3 seeds.

\begin{table}[t]
\centering
\small
\begin{tabular}{lccccc}
\toprule
\textbf{Condition} & \textbf{Seed 42} & \textbf{Seed 123} & \textbf{Seed 456} & \textbf{Mean} & \textbf{Std} \\
\midrule
Original & 99.83\% & 99.83\% & 99.83\% & 99.83\% & 0.00 \\
No-marker & 99.83\% & 99.83\% & 99.83\% & 99.83\% & 0.00 \\
\midrule
$\Delta$ & 0.00pp & 0.00pp & 0.00pp & 0.00pp & --- \\
\bottomrule
\end{tabular}
\caption{Format marker ablation (Qwen exception\_only). Removing the \texttt{[No exception data in trace]} marker produces zero degradation across all seeds (McNemar $p$=1.0).}
\label{tab:nomarker}
\end{table}

Table~\ref{tab:nomarker} shows zero degradation. McNemar's test yields $p$=1.0 (no difference). The error samples are identical across conditions (id=901, 905). This conclusively demonstrates the model learns exception fault semantics---the presence and content of error events---rather than format-level pattern matching.

\paragraph{Evidence Chain.} Table~\ref{tab:evidence_chain} summarizes the full evidence chain ruling out format shortcuts.

\begin{table}[t]
\centering
\small
\begin{tabular}{lcc}
\toprule
\textbf{Evidence} & \textbf{Value} & \textbf{Interpretation} \\
\midrule
Regex baseline (marker only) & 99.33\% & Marker \emph{is} a strong feature \\
SFT original (with marker) & 99.83\% & SFT $>$ regex by 0.50pp \\
SFT no-marker & 99.83\% & Removing marker $\rightarrow$ no loss \\
SFT $-$ Regex gap & +0.50pp & SFT learns beyond format \\
No-marker $-$ Original & 0.00pp & That learning $\neq$ the marker \\
Irreducible errors & 2/1190 & Same GeneratorExit samples \\
\bottomrule
\end{tabular}
\caption{Format shortcut falsification evidence chain. The model's 0.50pp advantage over regex comes from semantic understanding of exception content, not the marker string.}
\label{tab:evidence_chain}
\end{table}

\subsection{Claim 2: Information Density as Primary Bottleneck}

Full traces improve over no-trace (81.76\% vs 75.13\%, +6.6pp), but VOTA achieves 99.41\% (+24.3pp over no-trace). Label-only loss masking (oracle gradient on verdict token only) yields 86.22\% on full traces---quantifying gradient dilution at $\sim$26\% of the gap (4.5pp) and information overload at $\sim$74\% (13.0pp). Even with perfect loss targeting, full traces still underperform VOTA by 13pp, confirming information density as the dominant factor.

Cross-model evidence reinforces this: the VOTA gap between Qwen (99.41\%) and DeepSeek (84.45\%) is 14.96pp, partially attributable to training-time trace truncation (91.7\% of DeepSeek VOTA samples truncated from mean 2,626 to $\sim$2,048 tokens). Group analysis confirms a causal contribution: test samples whose training traces were truncated achieve 83.04\% vs 100\% for intact samples (Fisher exact $p \approx 0$).

\subsection{Claim 3: Two-Tier Diagnostic Hierarchy}

Component ablation on Qwen reveals a strict hierarchy:

\begin{itemize}
    \item \textbf{Tier 1 (Diagnostic Moments)}: exception\_only 99.83\%, ae\_only 99.92\%---individually near-perfect and sufficient ($p$=0.80 for Tier~1 internal comparison).
    \item \textbf{Tier 2 (Behavioral Patterns)}: loop\_only 85.63\%---substantial signal but model-dependent and fully subsumed by Tier~1.
\end{itemize}

Combining tiers introduces noise: exception\_only (99.83\%) $>$ VOTA full (99.41\%), $p$=0.01. The loop component approximates what label-only loss masking extracts from full traces (85.63\% $\approx$ 86.22\%), confirming VOTA's loop compression precisely captures the useful behavioral signal.

Calibration provides independent confirmation: Tier~1 ECE $\leq$ 0.003 versus no\_trace ECE = 0.172 ($\sim$63$\times$ gap); Brier scores show a $\sim$150$\times$ gap (0.0013 vs 0.1972). Exception signal not only improves accuracy but fundamentally improves model confidence calibration.

\subsection{Cross-Model Analysis}

\begin{table}[t]
\centering
\small
\begin{tabular}{lcccc}
\toprule
\textbf{Condition} & \textbf{Qwen-7B} & \textbf{DS-6.7B} & \textbf{CL-13B} & \textbf{Gap} \\
\midrule
exception\_only & 99.83\% & 99.33\%$^\dagger$ & 99.83\% & 0.50pp \\
VOTA & 99.41\% & 84.45\% & 92.94\% & 14.96pp \\
no\_trace & 75.13\% & 73.36\% & --- & 1.77pp \\
\bottomrule
\end{tabular}
\caption{Cross-model comparison. $^\dagger$DeepSeek median of 3 seeds (seed42/456=99.33\%, seed123=91.09\%). Gap = max cross-model difference. Exception\_only shows minimal gap; VOTA shows large model-dependent gap.}
\label{tab:cross_model}
\end{table}

Table~\ref{tab:cross_model} reveals that cross-model gap correlates with signal complexity. Exception\_only (structurally simplest) shows minimal gap (0.50pp between Qwen and DeepSeek median). VOTA (requires execution reasoning for Tier~2 loop component) shows 14.96pp gap. This confirms that simpler diagnostic signals generalize more robustly across architectures.

\paragraph{DeepSeek Seed Analysis.} DeepSeek exception\_only shows seed sensitivity: seed42/456 achieve 99.33\% (matching Qwen-level after accounting for 2 additional irreducible errors), while seed123 achieves 91.09\% (fail\_recall=0.705). Crucially, train\_loss is nearly identical across seeds (0.1225 vs 0.122), yet accuracy differs by 8.2pp---demonstrating that token-level loss masks sample-level classification ability in class-imbalanced settings. Even the worst seed (91.09\%) outperforms VOTA (84.45\%) and no\_trace (73.36\%), confirming relative dominance is robust.

\subsection{Irreducible Error Analysis}

The 2 errors (id=901, 905) consistent across all seeds and conditions originate from HumanEval/95 (\texttt{check\_dict\_case}). Root cause: Python's \texttt{all()} short-circuits on generator expressions, triggering \texttt{GeneratorExit} (a \texttt{BaseException} subclass used for control flow, not indicating actual errors). The tracer marks these as \texttt{[ERROR]}, causing the model to see 7 lines of \texttt{[ERROR] GeneratorExit} and reasonably predict fail. With only 2 such samples in the entire test set and no training signal for their benign nature, 99.83\% represents the theoretical ceiling for this representation.

\section{Discussion}

\paragraph{VOTA Repositioning.} Our results reposition VOTA from a generic trace compression method to a \textbf{fault indicator isolation} method. Its value lies not in making traces shorter, but in separating high-specificity fault indicators from low-specificity behavioral patterns. Exception\_only $\geq$ VOTA full demonstrates that the non-exception components are redundant for verdict prediction.

\paragraph{Training Stability as Secondary Factor.} DeepSeek's seed sensitivity (8.4pp gap between best and worst seed, with near-identical train loss) reveals LoRA initialization as a practical bottleneck for deploying exception-based verdict prediction on weaker models. We recommend multi-seed validation and training stabilization techniques (warmup scheduling, LoRA rank tuning) for production deployment.

\paragraph{Token-Level Loss Decoupling.} Both eval\_loss and train\_loss decouple from verdict accuracy: DeepSeek VOTA eval\_loss increases 65\% (epoch 1$\rightarrow$3) while accuracy improves 6.3pp; train\_loss convergence does not guarantee task performance across seeds. Verdict tasks require task-specific metrics for checkpoint selection.

\paragraph{Limitations.} (1) Our evaluation is limited to test verdict prediction on HumanEval + MBPP-sanitized. (2) DeepSeek VOTA gap is confounded by training-time truncation---the two factors (model capability, truncation) cannot be fully disentangled. (3) The 2,048-token training limit may disadvantage longer trace representations disproportionately.

\section{Related Work}

\paragraph{Trace-Augmented Code Models.} DCSH \cite{ni2024next} evaluated five trace formats for code generation, finding inconsistent benefits. TracePile \cite{jain2024tracepile} scaled trace SFT to 2.6M samples for generation tasks. TRACED \cite{dinella2022traced} used traces for encoder-based fault localization. Our work is the first systematic component ablation for verdict prediction.

\paragraph{Verdict Prediction.} Prior approaches use static features \cite{wang2021codet} or prompt-based reasoning. We show that minimal execution traces (exception events only) can achieve near-perfect accuracy, establishing a strong baseline for future work.

\section{Conclusion}

We demonstrate that exception events are consistently the most effective trace component for verdict prediction across 4 models, 3 architecture families, and 2 scales. The finding is robust to format-marker removal (zero degradation, McNemar $p$=1.0) and holds across all model families tested (exception dominance ranges from +0.25pp to +14.88pp over VOTA). A two-tier diagnostic hierarchy establishes that verdict prediction is fundamentally fault detection---the design principle is to extract highest-specificity diagnostic signals and discard everything else.

\bibliographystyle{plainnat}
\begin{thebibliography}{9}
\bibitem[Ni et al.(2024)]{ni2024next}
Ni, A., et al. (2024). Next-step execution trace prediction. \emph{EMNLP Findings}.

\bibitem[Jain et al.(2025)]{jain2024tracepile}
Jain, N., et al. (2025). TracePile: Scaling trace-augmented pretraining. \emph{NeurIPS}.

\bibitem[Dinella et al.(2024)]{dinella2022traced}
Dinella, E., et al. (2024). TRACED: Execution-aware pre-training for source code. \emph{ICSE}.

\bibitem[Wang et al.(2023)]{wang2021codet}
Wang, S., et al. (2023). CodeT: Code generation with generated tests. \emph{ICLR}.
\end{thebibliography}

\appendix

\section{CodeLlama-13B Detailed Results}
\label{app:codellama}

\begin{table}[h]
\centering
\small
\begin{tabular}{lc}
\toprule
\textbf{Metric} & \textbf{Value} \\
\midrule
Accuracy & 92.94\% \\
95\% CI & [91.4, 94.4] \\
$n$ & 1190 \\
Pass Precision & 0.93 \\
Pass Recall & 0.98 \\
Pass F1 & 0.955 \\
Fail Precision & 0.93 \\
Fail Recall & 0.80 \\
Fail F1 & 0.86 \\
Train Loss & 0.058 \\
Eval Loss & 0.12 \\
\bottomrule
\end{tabular}
\caption{CodeLlama-13B-Instruct VOTA performance. Lower fail recall (0.80) compared to Qwen (1.0) reflects weaker execution reasoning for Tier~2 loop components.}
\label{tab:codellama_detail}
\end{table}

\section{No-Marker Ablation Full Results}
\label{app:nomarker}

All three no-marker seeds produce identical results: accuracy 99.83\%, 2 errors (id=901, 905), pass\_precision=1.0, pass\_recall=0.9976, fail\_precision=0.9941, fail\_recall=1.0. Training losses: seed42=0.112, seed123=0.110, seed456=0.109. Eval losses: seed42=0.48, seed123=0.47, seed456=0.46. The monotonic decrease in loss across seeds is within normal variance and does not affect classification outcomes.

\end{document}
